% Resumen: Cinética química I (QUIM300). Compilar con pdflatex dos veces; usa resumen.sty.
\documentclass[10pt,letterpaper]{article}
\usepackage{resumen}

\newcommand{\A}{[\mathrm{A}]}
\newcommand{\B}{[\mathrm{B}]}
\newcommand{\vel}{\mathrm{M\,s^{-1}}}

% Más aire: separaciones locales (la plantilla no se modifica)
\setlength{\parskip}{5pt plus 1pt}
\renewcommand{\arraystretch}{1.3}
\titlespacing*{\section}{0pt}{15pt plus 3pt}{7pt}
\tcbset{
  clave/.append style={top=4pt,bottom=4pt,before skip=9pt,after skip=9pt},
  apoyo/.append style={top=4pt,bottom=4pt,before skip=8pt,after skip=8pt},
  nota/.append style={top=4pt,bottom=4pt,before skip=9pt,after skip=5pt},
  alerta/.append style={top=4pt,bottom=4pt,before skip=11pt,after skip=5pt}
}
\setlist{itemsep=4pt}

\begin{document}
\titulo{Cinética química I}{Velocidad de reacción y ley de velocidad}

%------------------------------------------------------------------
\section{Velocidad de reacción}
Para una reacción general $a\,\mathrm{A}+b\,\mathrm{B}\to c\,\mathrm{C}+d\,\mathrm{D}$, la velocidad de reacción se define
dividiendo el cambio de cada concentración por su coeficiente estequiométrico, de modo que una sola magnitud describe a
toda la reacción:

\begin{clave}
$\displaystyle \bar r=-\frac{1}{a}\frac{\Delta[\mathrm{A}]}{\Delta t}=-\frac{1}{b}\frac{\Delta[\mathrm{B}]}{\Delta t}
=\frac{1}{c}\frac{\Delta[\mathrm{C}]}{\Delta t}=\frac{1}{d}\frac{\Delta[\mathrm{D}]}{\Delta t}$
\hfill\small velocidad media en $\Delta t$
\end{clave}
\begin{clave}
$\displaystyle r=\lim_{\Delta t\to0}\bar r=-\frac{1}{a}\frac{d[\mathrm{A}]}{dt}=-\frac{1}{b}\frac{d[\mathrm{B}]}{dt}
=\frac{1}{c}\frac{d[\mathrm{C}]}{dt}=\frac{1}{d}\frac{d[\mathrm{D}]}{dt}$
\hfill\small velocidad instantánea
\end{clave}

\vspace{4pt}
\begin{minipage}[t]{0.48\linewidth}\vspace{0pt}
Como los reactivos disminuyen ($\Delta[\mathrm{R}]<0$), el signo menos hace que $r$ sea positiva. La velocidad $r$ es
única para la reacción completa, sea cual sea la especie con que se mida, y sus unidades son concentración sobre tiempo,
$\mathrm{mol\,L^{-1}\,s^{-1}}$ ($\vel$).\par\vspace{6pt}
En el gráfico de $\A$ contra $t$, la velocidad media es el negativo de la pendiente de una secante y la instantánea, el
negativo de la pendiente de la tangente.

\begin{apoyo}
\textbf{Ejemplo.} $\A$ baja de \SI{0,80}{\molar} a \SI{0,50}{\molar} en \SI{60}{\second}:\\[4pt]
\centering $\displaystyle \bar r=-\frac{(0{,}50-0{,}80)\ \mathrm{M}}{60\ \mathrm{s}}=\SI{5,0e-3}{\molar\per\second}$
\end{apoyo}
\end{minipage}\hfill
\begin{minipage}[t]{0.46\linewidth}\vspace{0pt}
\begin{tikzpicture}
\begin{axis}[resumen,width=\linewidth,height=4.9cm,xlabel={$t$ (s)},ylabel={$[\mathrm{A}]$ (M)},
  xmin=0,xmax=300,ymin=0,ymax=1.05,xtick={0,100,200,300},ytick={0,0.5,1},clip=false]
\addplot[acc,domain=0:300,samples=60,line width=1.5pt]{exp(-x/120)};
% secante entre t=100 y t=240
\addplot[muted,line width=1pt,domain=100:240]{0.43460-0.0021381*(x-100)};
\addplot[muted,only marks,mark=*,mark size=1.6pt] coordinates {(100,0.43460)(240,0.13534)};
\draw[muted,dashed] (axis cs:100,0.13534)--(axis cs:240,0.13534)--(axis cs:240,0.13534);
\draw[muted,dashed] (axis cs:100,0.13534)--(axis cs:100,0.43460);
\node[font=\scriptsize,muted] at (axis cs:170,0.07) {$\Delta t$};
\node[font=\scriptsize,muted,anchor=east] at (axis cs:96,0.28) {$\Delta[\mathrm{A}]$};
% tangente en t=40
\addplot[bad,line width=1pt,domain=0:92]{0.71653-0.0059711*(x-40)};
\addplot[bad,only marks,mark=*,mark size=2pt] coordinates {(40,0.71653)};
\node[font=\scriptsize,bad,anchor=west] at (axis cs:55,0.93) {tangente: pendiente $=-r$};
\node[font=\scriptsize,muted,anchor=west] at (axis cs:150,0.60) {secante: pendiente $=-\bar r$};
\end{axis}
\end{tikzpicture}
\end{minipage}

%------------------------------------------------------------------
\section{Aparición y consumo}
Cada especie cambia su concentración con una rapidez proporcional a su coeficiente estequiométrico $|\nu_i|$. Por eso las
velocidades de consumo o de aparición difieren entre especies, mientras que $r$ es única:

\vspace{2pt}
\begin{minipage}[t]{0.38\linewidth}\vspace{0pt}
\begin{clave}
$\displaystyle r_i=|\nu_i|\,r$
\end{clave}
\small $-\dfrac{d[\mathrm{R}]}{dt}$: consumo de un reactivo.\\[5pt]
$+\dfrac{d[\mathrm{P}]}{dt}$: aparición de un producto.
\end{minipage}\hfill
\begin{minipage}[t]{0.58\linewidth}\vspace{0pt}
\small\setlength{\tabcolsep}{6pt}
\begin{tabular}{@{}lccc@{}}
\toprule
$\mathrm{N_2}+3\,\mathrm{H_2}\to2\,\mathrm{NH_3}$ & $\mathrm{N_2}$ & $\mathrm{H_2}$ & $\mathrm{NH_3}$\\
\midrule
Tipo & consumo & consumo & aparición\\
$r_i$ & $1\,r$ & $3\,r$ & $2\,r$\\
valor ($10^{-3}\ \mathrm{M\,s^{-1}}$) & $1{,}0$ & $3{,}0$ & $2{,}0$\\
\bottomrule
\end{tabular}\\[3pt]
{\footnotesize Valores para $r=1{,}0\times10^{-3}\ \mathrm{M\,s^{-1}}$.}
\end{minipage}

%------------------------------------------------------------------
\section{Ley de velocidad y órdenes de reacción}
A temperatura constante, la velocidad depende de las concentraciones de los reactivos:
\begin{clave}
$\displaystyle r=k\,\A^{m}\,\B^{n}$\hfill\small orden total $=m+n$
\end{clave}

\vspace{2pt}
\begin{minipage}[t]{0.47\linewidth}\vspace{0pt}
La constante de velocidad $k$ depende de la reacción, de la temperatura y del catalizador, pero no de las
concentraciones. Los órdenes parciales $m$ y $n$ se determinan experimentalmente.\par\vspace{6pt}
Si $\A$ se multiplica por $f$ con $\B$ constante, $r$ se multiplica por $f^{m}$:
\begin{apoyo}
$\displaystyle \frac{r_2}{r_1}=\frac{k\,(f\A)^{m}}{k\,\A^{m}}=f^{m}$
\end{apoyo}
\end{minipage}\hfill
\begin{minipage}[t]{0.49\linewidth}\vspace{0pt}
\begin{tikzpicture}
\begin{axis}[resumen,width=0.86\linewidth,height=4.6cm,xlabel={$[\mathrm{A}]/[\mathrm{A}]_0$},ylabel={$r/r_0$},
  xmin=0,xmax=3,ymin=0,ymax=9,xtick={0,1,2,3},ytick={0,3,6,9},clip=false]
\addplot[muted,dash pattern=on 2pt off 2pt,domain=0:3]{1};
\addplot[teal,dashed,domain=0:3]{x};
\addplot[acc,domain=0:3,samples=60]{x^2};
\node[font=\scriptsize,muted,anchor=west] at (axis cs:3.05,1) {orden 0};
\node[font=\scriptsize,teal,anchor=west] at (axis cs:3.05,3.2) {orden 1};
\node[font=\scriptsize,acc,anchor=west] at (axis cs:3.05,8.6) {orden 2};
\end{axis}
\end{tikzpicture}

\vspace{12pt}
\small\textbf{Unidades de $k$.} Con orden total $N$ y concentraciones en M, $k$ queda en $\mathrm{M^{1-N}\,s^{-1}}$:\\[6pt]
\begin{tabular}{@{}lcccc@{}}
\toprule
$N$ & $0$ & $1$ & $2$ & $3$\\
\midrule
$k$ & $\mathrm{M\,s^{-1}}$ & $\mathrm{s^{-1}}$ & $\mathrm{M^{-1}\,s^{-1}}$ & $\mathrm{M^{-2}\,s^{-1}}$\\
\bottomrule
\end{tabular}
\end{minipage}

\begin{alerta}
\textbf{Errores frecuentes.} Los órdenes no se deducen de la ecuación balanceada: solo coinciden con los coeficientes en un
paso elemental. Por ejemplo, $2\,\mathrm{NO}+2\,\mathrm{H_2}\to\mathrm{N_2}+2\,\mathrm{H_2O}$ cumple
$r=k[\mathrm{NO}]^{2}[\mathrm{H_2}]$, con orden 1 en $\mathrm{H_2}$ aunque su coeficiente es 2, y orden total 3.
\end{alerta}

%------------------------------------------------------------------
\section{Método de las velocidades iniciales}
Permite hallar los órdenes parciales y $k$ a partir de una serie de experimentos a la misma temperatura. La velocidad
se mide al comienzo de cada experimento, cuando aún no hay productos que interfieran.

\vspace{2pt}
\begin{minipage}[t]{0.47\linewidth}\vspace{0pt}
\begin{enumerate}[label=\textbf{\arabic*.}]
\item Se mide la velocidad inicial $r_0$ con distintas concentraciones iniciales.
\item Se eligen dos experimentos en los que cambia la concentración de una sola especie.
\item Se comparan las velocidades:
\end{enumerate}
\begin{clave}
$\displaystyle \frac{r_2}{r_1}=\left(\frac{\A_2}{\A_1}\right)^{m}$
\end{clave}
\begin{enumerate}[label=\textbf{\arabic*.},start=4]
\item Se repite el procedimiento para cada reactivo, se suman los órdenes parciales y se calcula $k=r_0/(\A^{m}\B^{n})$, que
se verifica con otro experimento.
\end{enumerate}
\begin{nota}
\textbf{Verificación de sanidad.} El orden obtenido debe concordar con lo esperado al duplicar $\A$: si $r$ no cambia,
$m=0$; si se duplica, $m=1$; si se cuadruplica, $m=2$; si se multiplica por 8, $m=3$.
\end{nota}
\vspace{6pt}
En el ejemplo, la ley de velocidad resulta $r=k\,[\mathrm{NO}]^{2}[\mathrm{H_2}]$, de orden total 3.
\end{minipage}\hfill
\begin{minipage}[t]{0.49\linewidth}\vspace{0pt}
\textbf{Ejemplo:} $2\,\mathrm{NO}+2\,\mathrm{H_2}\to\mathrm{N_2}+2\,\mathrm{H_2O}$\\[5pt]
\small
\begin{tabular}{@{}cccc@{}}
\toprule
Exp. & $[\mathrm{NO}]$ (M) & $[\mathrm{H_2}]$ (M) & $r_0$ ($\mathrm{M\,s^{-1}}$)\\
\midrule
1 & 0,010 & 0,010 & \num{2,5e-4}\\
2 & 0,020 & 0,010 & \num{1,0e-3}\\
3 & 0,010 & 0,020 & \num{5,0e-4}\\
\bottomrule
\end{tabular}\\[8pt]
\normalsize
Exp.\ 1 y 2 (cambia $[\mathrm{NO}]$):\vspace{-2pt}
\begin{apoyo}\small
$\displaystyle\frac{r_2}{r_1}=\frac{1{,}0\times10^{-3}}{2{,}5\times10^{-4}}=4=\left(\frac{0{,}020}{0{,}010}\right)^{m}=2^{m}\Rightarrow m=2$
\end{apoyo}
Exp.\ 1 y 3 (cambia $[\mathrm{H_2}]$):\vspace{-2pt}
\begin{apoyo}\small
$\displaystyle\frac{r_3}{r_1}=\frac{5{,}0\times10^{-4}}{2{,}5\times10^{-4}}=2=\left(\frac{0{,}020}{0{,}010}\right)^{n}=2^{n}\Rightarrow n=1$
\end{apoyo}
Constante, con el Exp.\ 1:\vspace{-2pt}
\begin{clave}\small
$\displaystyle\begin{aligned}k&=\frac{r_0}{[\mathrm{NO}]^{2}[\mathrm{H_2}]}=\frac{2{,}5\times10^{-4}\ \mathrm{M\,s^{-1}}}{(0{,}010\ \mathrm{M})^{2}(0{,}010\ \mathrm{M})}\\
&=2{,}5\times10^{2}\ \mathrm{M^{-2}\,s^{-1}}\end{aligned}$
\end{clave}
\end{minipage}

\end{document}
